\documentclass[10pt,a4paper]{amsart}
\usepackage[margin=19mm]{geometry}
\usepackage{amsmath,amssymb,mathtools}
\usepackage[scr=rsfs,cal=euler]{mathalfa}
\usepackage{xfrac}
\usepackage[unicode,hidelinks,bookmarks=false]{hyperref}
\newcommand{\ie}{i.e.\ }
\newcommand{\cf}{cf.\ }
\newcommand{\resp}{respectively\ }
\newcommand{\Subsection}{\subsection}
\providecommand{\nobreakdash}{\nobreak\hbox{-}}
\setlength{\emergencystretch}{2em}
\multlinegap0pt
\usepackage{bm}
\usepackage{bbm}
\usepackage{dsfont}
\usepackage{xspace}
\usepackage{cancel}
\usepackage{mathtools}
\usepackage{amsthm}
\makeatletter
\@ifundefined{lemma}{\newtheorem{lemma}{Lemma}}{}
\@ifundefined{prop}{\newtheorem{prop}{Proposition}}{}
\makeatother
\newcommand{\RR}{\mathbb{R}}
\newcommand{\ZZ}{\mathbb{Z}}
\title[Valuatively independent bases for the Fermat family of cubic]{Valuatively independent bases for the Fermat family of cubic curves}
\author{Jakob Hultgren, Sohaib Khalid}
\date{Source: arXiv:2604.16296v1 (CC BY 4.0)}
\begin{document}
\maketitle

\noindent\emph{Source and licence.} Valuatively independent bases for the Fermat family of cubic curves. Version arXiv:2604.16296v1 (CC BY 4.0), distributed under the Creative Commons Attribution 4.0 International licence, as recorded in the arXiv licence metadata for this exact version and shown on the abstract page https://arxiv.org/abs/2604.16296v1. Full rights notice: licenses/arxiv-2604.16296-CC-BY-4.0.txt.\par\smallskip

\noindent\textit{Editorial scope.} Counted unit: the Prop whose source label is \texttt{None} in the article (source0.tex lines 206--215, source label \texttt{None}), delivered together with the article's own complete original proof of it (source0.tex lines 217--268, one \texttt{proof} environment of 52 lines). No mathematics is written, repaired or re-derived: statement and proof are the article's own text line for line. Required context retained uncounted: lemma:AffineIso (lines 152--162). Every definition, display and label of those lines is retained unchanged. Omitted from this package and not redistributed: 1--151, 163--205, 216--216, 269--756. Nothing inside a retained excerpt is omitted or altered: every retained body is delivered exactly as the article's lines are written, so no omission receipt of any kind accompanies this package. Citations: the retained excerpts carry no citation command at all, so no bibliography entry of the article is transcribed and no reference is redistributed. Reference adaptation: none was needed and none was made; every reference inside the delivered unit resolves inside it, so no unresolved reference or citation is produced. Printed slips kept literal: no printed word, symbol, hypothesis or number of the counted statement or of its proof was altered. Layout and class: the article's own class is \texttt{article} with packages including amsfonts, amsmath, amssymb, amsthm, tikz, graphicx, geometry; the wrapper is the lane's pinned \texttt{amsart} editorial wrapper at 10pt on A4, which restates the theorem-like environments the retained text needs and adds the article's own macro definitions that the retained text needs (\texttt{\textbackslash RR}, \texttt{\textbackslash ZZ}), with extra wrapper packages (bm, bbm, dsfont, xspace, cancel, mathtools). The delivered numbering is the unit's own. Assets: no retained span contains a figure environment, a graphics-inclusion command, a PDF-page inclusion, a table or a TikZ picture; no image file is copied into this package, the package declares no asset dependency and no image file is redistributed with it. Licence: version arXiv:2604.16296v1 of this article is distributed under the Creative Commons Attribution 4.0 International licence, as recorded in the arXiv licence metadata for this exact version and shown on the abstract page https://arxiv.org/abs/2604.16296v1. A full rights notice is delivered at licenses/arxiv-2604.16296-CC-BY-4.0.txt. Characters: every retained excerpt is pure ASCII, so the delivered engine, XeLaTeX with its default Latin Modern text font, reports no missing character.
\begin{lemma}\label{lemma:AffineIso}
    There are affine isomorphisms 
    \[
L \cong \frac{\RR\times\RR}{\langle (t,a)\mapsto (t-3,a-9t+9)\rangle} \longrightarrow M \cong \frac{\RR}{\langle t\mapsto t-3\rangle},
\]
and 
\[
L^\star \cong \frac{\RR\times\RR}{\langle (t^\vee,b)\mapsto (t^\vee-9,b+3t^\vee-18)\rangle} \longrightarrow M^\star \cong \frac{\RR}{\langle t^\vee\mapsto t^\vee-9\rangle}
\]
which are compatible with the action of $\Pi\cong\ZZ$ on the universal covers  of $M,M^\star,L,L^\star$.
\end{lemma}

\begin{prop}
    In terms of the affine isomorphisms furnished by Lemma~\ref{lemma:AffineIso}, the cost function $c:M^\star\times M \to \RR$ is (induced from the map $\RR\times\RR \to \RR$) given by
    \[
c(t^\vee,t) = 3(k+1)t + \ell t^\vee-(t^\vee -9m)(t+3m) - \frac{3k(k+1)+3\ell(\ell+1)+9m(3m-1)}{2}
\]
where
\[
k= \lfloor t \rfloor, \quad \ell = \left\lfloor \frac{t^\vee}{3}\right\rfloor, \quad m =\left \lfloor \frac{t^\vee-3t}{9}+\frac{1}{3}\right\rfloor = \left\lfloor \frac{1}{3}\left(\frac{t^\vee}{3}-t + 1 \right)\right\rfloor.
\]
\end{prop}

\begin{proof}
Under the isomorphisms furnished by the Lemma, the pullback $\pi^*\Phi_0$ of the piecewise affine convex section $\Phi_0$ of $L$  is given by 
\[
\pi^*\Phi_0(t)=\left(t,\sup_{k\in\ZZ}\left(3kt-1-\frac{3k(k-1)}{2}\right)\right) = \left(t,3(\lfloor t \rfloor +1)t-1-\frac{3\lfloor t \rfloor(\lfloor t \rfloor+1)}{2}\right).
\]
The section of $[\cdot,\cdot]$ of $L\boxplus-L^\star \to M\times M^\star$ under this isomorphism is induced by the pairing
\[
[t,t^\vee] = \sup_{k\in\ZZ}(\gamma^k \cdot s_{(t^\vee,0)})(t) - s_{(t^\vee,0)}.
\]
(Note that we could have used $s_{(t^\vee,b)}$ for any $b\in \RR$ in the above expression without affecting its value, since the definition of the $\RR$-action on $L\boxplus-L^\star$ means that $(t,a)-s_{(t^\vee,b)}$ and $(t,a-b)-s_{(t^\vee,0)}$ represent the same points in the fibre over $(t,t^\vee)$.) In order to write this explicitly, we need to determine the supremum
\begin{align*}
    \sup_{k\in\ZZ}(\gamma^k \cdot s_{(t^\vee,0)})(t) &= \left(t,t^\vee t +\sup_{k\in\ZZ}\left(3kt^\vee-9kt-18k-27\frac{k(k-1)}{2}\right)\right)\\
    &=\left(t,\sup_{k\in\ZZ} \left((t^\vee-9k)(t+3k)+27k^2-18k-27\frac{k(k-1)}{2}\right)\right)
\end{align*}
for any given value of $(t,t^\vee)$. Now 

\begin{align*}
    3kt^\vee-9kt-18k-27\frac{k(k-1)}{2}&=\frac{9k}{2}\left(2(t^\vee/3 - t)-4-3(k-1)\right)\\
    &=\frac{9k}{2}\left(2(t^\vee/3 - t)-1-3k\right)\\
    &= \frac{9k}{2}(6r-3k)\\
    &= \frac{27}{2}(r^2-(k-r)^2).
\end{align*}
where for convenience we have written $r=\frac{1}{9}(t^\vee-3t)-\frac{1}{6}$. This last expression attains its largest value when $|k-r|$ attains its smallest value. If $\{r\} = r-\lfloor r \rfloor \leq\frac{1}{2}$, then this happens when $k=\lfloor r \rfloor$ and if $\{r\} \geq \frac{1}{2}$, this happens for $k=\lfloor r \rfloor +1$. Thus, we get 
\[
[t,t^\vee] =
\begin{cases}
    (t, t^\vee t + \frac{27}{2}(r^2-\{r\}^2))-s_{(t^\vee,0)}, & \textrm{ if } \{r\}\leq \frac{1}{2},\\
    \\
    (t, t^\vee t + \frac{27}{2}(r^2-(1-\{r\})^2))-s_{(t^\vee,0)}, & \textrm{ if } \{r\}\geq \frac{1}{2}.
\end{cases}
\]
Note that the two expressions clearly coincide when $\{r\} = \frac{1}{2}$. When $r\in[k-1/2,k+1/2]$ for $k\in \ZZ$, we can easily check that the two expressions both reduce to $t^\vee t + \frac{27}{2}(r^2-(r-k)^2)$ (with both boundary points giving the same value $\frac{27}{2}(r^2+\frac{1}{4})$). Now $r\in [k-1/2,k+1/2)$ if and only if $k=\left\lfloor \frac{t^\vee-3t}{9}+\frac{1}{3}\right\rfloor$ and so we obtain
\[
[t,t^\vee]=\left(t,(t^\vee-9k)(t+3k)+\frac{9k(3k-1)}{2}\right) -s_{(t^\vee,0)} \textrm{ where } k=\left\lfloor \frac{t^\vee-3t}{9}+\frac{1}{3}\right\rfloor.
\]
Given a continuous section $\Phi$ of $L$, its \emph{Legendre transform} $\Phi^\star$ is a section (induced by the function) defined by the formula 
\[
\Psi^\star(t^\vee) = \sup_{t\in \RR} \left([t,t^\vee]-\Psi(t)\right).
\]
and is a continuous section of $-L^\star$ by construction. We shall use $\Phi_0$ and $\Psi = \Phi_0^\star$ to write down the expression for $c$, where the section $\Phi_0$ is the convex section which on $U_i$ is given by $n\mapsto -\langle m_i,n\rangle$. We can write down explicitly:
\[
\Phi_0^\star(t^\vee)=\sup_{t\in \RR} \left([t,t^\vee]-\Phi_0(t)\right) = \left(t^\vee, \lfloor t^\vee /3 \rfloor t + 1-\frac{3\lfloor t^\vee /3 \rfloor(\lfloor t^\vee /3 \rfloor+1)}{2} \right).
\]
Finally, the cost function $c:M\times M^\star\to \RR$ is induced by the function (which we still denote by $c$) given by $c(t,t^\vee)=-[t,t^\vee]+\Phi_0(t)+\Phi_0^\star(t^\vee)$ which we can write explicitly as 
\begin{equation}\label{eq:Cost}
c(t,t^\vee) = 3(k+1)t + \ell t^\vee-(t^\vee -9m)(t+3m) - \frac{3k(k+1)+3\ell(\ell+1)+9m(3m-1)}{2}
\end{equation}
where
\[
k= \lfloor t \rfloor, \quad \ell = \left\lfloor \frac{t^\vee}{3}\right\rfloor, \quad m =\left \lfloor \frac{t^\vee-3t}{9}+\frac{1}{3}\right\rfloor.
\]
\end{proof}

\end{document}
